\section{Masalah Pengalamatan Memori dan Lahirnya Segmentasi}
\label{sec:masalah-segmentasi}

Program assembly bekerja langsung dengan alamat memori. Pada
Bab~\ref{chap:register} telah dibahas bahwa 8086 memiliki empat
register segmen, yaitu \reg{CS}, \reg{DS}, \reg{SS}, dan \reg{ES},
yang masing-masing berukuran 16 bit. Muncul pertanyaan yang wajar:
mengapa prosesor yang bus alamatnya berukuran 20 bit justru memakai
register segmen 16 bit? Jawaban atas pertanyaan itulah yang menjadi
inti bab ini.

\subsection{Keterbatasan Register 16-Bit}

Sebuah register 16-bit hanya mampu menampung nilai dari \hex{0000}
sampai \hex{FFFF}, yaitu \dec{65.536} nilai yang berbeda. Bila satu
register dipakai untuk menunjuk satu alamat memori secara langsung,
maka jangkauan alamat yang dapat dijangkau hanyalah \dec{65.536}
byte atau 64 KB. Dengan kata lain, seluruh memori yang dapat
disentuh program akan terbatas pada wilayah 64 KB saja.

Keadaan menjadi berbeda ketika kita melihat sisi perangkat keras.
Prosesor 8086 memiliki bus alamat 20 bit, sehingga ia mampu
menjangkau $2^{20} = \dec{1.048.576}$ alamat, yaitu 1 MB
(\eng{megabyte}) memori. Angka ini enam belas kali lebih besar
daripada jangkauan satu register 16-bit. Perbandingan tersebut
dirangkum pada Tabel~\ref{tab:lebar-alamat}.

\begin{table}[H]
    \centering
    \caption{Perbandingan lebar register dan lebar bus alamat 8086.}
    \label{tab:lebar-alamat}
    \small
    \begin{tabular}{@{}p{3.6cm}p{3.4cm}p{3.0cm}p{4.0cm}@{}}
        \toprule
        \textbf{Sisi} & \textbf{Lebar} & \textbf{Jangkauan} &
        \textbf{Akibat} \\
        \midrule
        Register umum dan register segmen & 16 bit &
        \dec{65.536} alamat & Hanya 64 KB dapat ditunjuk langsung \\
        \addlinespace
        Bus alamat prosesor & 20 bit &
        \dec{1.048.576} alamat & Seluruh 1 MB dapat dijangkau \\
        \bottomrule
    \end{tabular}
\end{table}

Terdapat selisih jangkauan sebesar 16 kali antara kedua sisi
tersebut. Jika register 16-bit dipakai seorang diri, memori di atas
alamat \hex{FFFF} tidak akan pernah dapat diakses, padahal peta
memori 8086 justru menempatkan bagian-bagian penting sistem pada
alamat-alamat tinggi, misalnya ROM BIOS pada \hex{F0000}. Di sinilah
sebuah strategi tambahan diperlukan, dan strategi itu bernama
\term{segmentasi}.

\begin{infobox}
Segmentasi bukan sekadar keterbatasan yang harus diterima, melainkan
sebuah kompromi rancangan. Dengan segmentasi, lebar instruksi dan
register tetap 16 bit (sehingga program tetap ringkas dan kompatibel
dengan perangkat lunak lama), tetapi jangkauan alamat tetap mencapai
1 MB. Konsep serupa dipakai pada hampir semua sistem berbasis 8086
dan diwarisi hingga prosesor modern sebagai mode real.
\end{infobox}

\subsection{Alamat Logis dan Alamat Fisik}

Pemecahan masalah di atas dilakukan dengan membagi sebuah alamat
menjadi dua bagian yang ditulis berdampingan dan dipisahkan tanda
titik dua. Bentuk penulisan seperti \code{1234h:5678h} disebut
\term{alamat logis}, sedangkan alamat sebenarnya yang dikirim ke
bus alamat disebut \term{alamat fisik}.

\begin{definisi}
\textbf{Alamat logis} adalah cara programmer menuliskan sebuah
alamat, yaitu berupa pasangan \eng{segment:offset} dengan bentuk
\code{segmen:offset}. \textbf{Alamat fisik} adalah alamat 20-bit
yang benar-benar dibangkitkan prosesor dan dikirim ke bus alamat
untuk memilih satu lokasi memori.
\end{definisi}

Pada penulisan \code{1234h:5678h}, nilai \hex{1234} disebut
\term{segmen} (\eng{segment}) dan nilai \hex{5678} disebut
\term{offset}. Keduanya berukuran 16 bit. Titik dua di antara
keduanya hanya berfungsi sebagai pemisah penulisan; prosesor tidak
mengenalnya sebagai operasi pembagian. Ketika programmer menulis
\lstasm|MOV AX, [1234h]|, yang dipakai adalah segmen data saat itu
(ditunjuk oleh \reg{DS}) dengan offset \hex{1234}. Sebaliknya, bila
alamat ditulis lengkap sebagai \lstasm|MOV AX, 1234h:[5678h]|,
pasangan tersebut dipakai apa adanya tanpa memperhatikan isi
\reg{DS}.

\begin{infobox}
\important{Nilai segmen bukan alamat.} Segmen \hex{1234} bukan
menunjuk alamat \hex{1234}, melainkan menunjuk awal sebuah paragraf
yang letaknya lebih jauh, yaitu pada \hex{12340}. Kekeliruan
memperlakukan nilai segmen sebagai alamat sudah menjadi sumber
kesalahan perhitungan yang paling sering terjadi.
\end{infobox}

\subsection{Rumus Alamat Fisik}

Hubungan antara alamat logis dan alamat fisik dinyatakan oleh rumus
berikut.

\begin{equation}
\label{eq:rumus-alamat-fisik}
\text{Alamat fisik} = (\text{Segmen} \times \hex{10}) +
\text{Offset}
\end{equation}

Nilai \hex{10} pada Persamaan~\ref{eq:rumus-alamat-fisik} adalah
bilangan heksadesimal yang sama dengan \dec{16} desimal. Prosesor
tidak benar-benar melakukan perkalian; secara perangkat keras ia
hanya menggeser nilai segmen empat bit ke kiri (setara dengan
mengalikan 16), lalu menjumlahkannya dengan offset. Karena segmen
dan offset masing-masing 16 bit, hasil penjumlahan tersebut dapat
mencapai 20 bit, bahkan pada kasus tertentu menghasilkan bit ke-21
yang akan dibahas pada Bagian~\ref{sec:peta-memori}.

Mengapa pengalinya justru 16 dan bukan 10 atau 100? Karena memori
dibagi menjadi satuan-satuan berukuran 16 byte yang disebut
\term{paragraf} (\eng{paragraph}). Nilai segmen sesungguhnya
menyatakan nomor paragraf, bukan nomor byte. Bila segmen bernilai
\hex{1234}, artinya alamat dasar berada pada paragraf ke-\hex{1234},
dan setiap paragraf berisi 16 byte. Dengan demikian dasar segmen
tersebut berada pada byte ke $(\hex{1234} \times \dec{16}) =
\hex{12340}$. Satuan inilah yang membuat nilai segmen selalu
berakhir pada digit nol ketika dituliskan sebagai alamat fisik.

Gambar~\ref{fig:pembentukan-alamat} memperlihatkan proses
pembentukan alamat fisik untuk pasangan \code{1234h:5678h}. Perhatikan
bahwa pada baris segmen ditambahkan satu digit nol di sebelah kanan,
yang menyatakan pergeseran empat bit tersebut.

\begin{figure}[H]
    \centering
    \begin{tikzpicture}[
        digit/.style={draw=obeBlue, line width=0.5pt, minimum size=0.65cm,
                      inner sep=0pt, font=\ttfamily\small, anchor=center},
        ket/.style={font=\scriptsize, align=center},
        labelkiri/.style={font=\scriptsize\bfseries, text=obeBlue,
                          anchor=east}
    ]
        % ---- baris segmen yang sudah digeser 4 bit ----
        \foreach \d [count=\i from 0] in {1,2,3,4}
            \node[digit, fill=obeBlueLight] at (1.0+0.65*\i, 2.2) {\d};
        \node[digit, fill=obeTealLight, draw=obeTeal] (nol) at (3.60, 2.2) {0};

        % ---- baris offset ----
        \foreach \d [count=\i from 1] in {5,6,7,8}
            \node[digit, fill=obeGrayLight, draw=obeGray]
                at (1.0+0.65*\i, 1.45) {\d};

        % ---- garis penjumlahan ----
        \draw[line width=0.9pt, obeBlue] (0.675, 1.02) -- (3.925, 1.02);

        % ---- baris hasil ----
        \foreach \d [count=\i from 0] in {1,7,9,B,8}
            \node[digit, fill=obeAmberLight, draw=obeAmber]
                at (1.0+0.65*\i, 0.45) {\d};

        % ---- keterangan di sebelah kiri ----
        \node[labelkiri] at (0.55, 2.2) {Segmen};
        \node[labelkiri] at (0.55, 1.45) {Offset};
        \node[labelkiri] at (0.55, 0.45) {Fisik};

        % ---- anotasi pergeseran ----
        \draw[-{Stealth[length=2mm]}, obeTeal, line width=0.7pt]
            (nol.north) -- ++(0,0.42)
            node[above, ket, text=obeTeal]
            {digit 0 ditambahkan:\\segmen digeser 4 bit ($\times 16$)};
    \end{tikzpicture}
    \caption{Pembentukan alamat fisik 20-bit dari pasangan
    \eng{segment:offset}. Segmen digeser empat bit, lalu dijumlahkan
    dengan offset.}
    \label{fig:pembentukan-alamat}
\end{figure}

\begin{contoh}
Hitunglah alamat fisik dari pasangan \code{1234h:5678h}.\\[0.3em]
Langkah pertama, kalikan segmen dengan \hex{10}:
$\hex{1234} \times \hex{10} = \hex{12340}$. Langkah kedua,
jumlahkan hasilnya dengan offset:
$\hex{12340} + \hex{5678} = \hex{179B8}$. Jadi alamat fisiknya
adalah \hex{179B8}.
\end{contoh}

\subsection{Contoh Perhitungan Bertingkat}

Enam contoh berikut disusun bertingkat, mulai dari kasus paling
sederhana sampai kasus yang memunculkan \eng{carry} pada digit
kelima. Seluruh hasilnya dirangkum pada
Tabel~\ref{tab:contoh-alamat-fisik}. Kerjakan setiap baris dengan
rumus pada Persamaan~\ref{eq:rumus-alamat-fisik} sebelum melihat
hasilnya. Pada kolom terakhir, alamat fisik ditulis dengan lima
digit heksadesimal agar lebar 20-bitnya terlihat; nilai \hex{00100}
sama dengan \hex{0100}.

\begin{table}[H]
    \centering
    \caption{Enam contoh perhitungan alamat fisik yang disusun
    bertingkat.}
    \label{tab:contoh-alamat-fisik}
    \small
    \begin{tabular}{@{}cp{1.6cm}p{1.6cm}p{3.9cm}p{2.1cm}@{}}
        \toprule
        \textbf{No} & \textbf{Segmen} & \textbf{Offset} &
        \textbf{Perhitungan} & \textbf{Alamat fisik} \\
        \midrule
        1 & \hex{0000} & \hex{0100} &
        \hex{00000} + \hex{0100} & \hex{00100} \\
        \addlinespace
        2 & \hex{1000} & \hex{0020} &
        \hex{10000} + \hex{0020} & \hex{10020} \\
        \addlinespace
        3 & \hex{076A} & \hex{0100} &
        \hex{076A0} + \hex{0100} & \hex{077A0} \\
        \addlinespace
        4 & \hex{2000} & \hex{FFFF} &
        \hex{20000} + \hex{FFFF} & \hex{2FFFF} \\
        \addlinespace
        5 & \hex{1234} & \hex{5678} &
        \hex{12340} + \hex{5678} & \hex{179B8} \\
        \addlinespace
        6 & \hex{B800} & \hex{0000} &
        \hex{B8000} + \hex{0000} & \hex{B8000} \\
        \bottomrule
    \end{tabular}
\end{table}

Setiap baris mengajarkan hal yang berbeda. Baris 1 memperlihatkan
segmen \hex{0000}, yaitu segmen paling awal, sehingga alamat
fisiknya sama dengan offsetnya sendiri. Baris 2 memperlihatkan
segmen yang berakhir pada nol sehingga penjumlahannya mudah dibaca.
Baris 3 adalah bentuk yang paling sering muncul dalam praktikum
karena segmennya diambil dari nilai \reg{DS} yang diberikan DOS.
Baris 4 memakai offset terbesar yang mungkin, yaitu \hex{FFFF},
sehingga menghasilkan alamat tertinggi yang dapat dicapai segmen
tersebut. Baris 5 memunculkan \eng{carry} yang mengubah digit
kelima hasil penjumlahan. Baris 6 menunjuk awal area video teks
warna, yang akan dibahas pada Bagian~\ref{sec:peta-memori}.

\begin{peringatanbox}
Kesalahan paling sering adalah menjumlahkan segmen dan offset tanpa
mengalikan segmen dengan \hex{10}. Pada contoh \code{1234h:5678h},
cara yang salah menghasilkan $\hex{1234} + \hex{5678} = \hex{68AC}$,
sedangkan hasil yang benar adalah \hex{179B8}. Selisih kedua hasil
tersebut sangat besar karena satu langkah perkalian 16 telah
terlewat. Biasakan menuliskan nilai segmen yang sudah digeser
sebelum melakukan penjumlahan.
\end{peringatanbox}

Dengan rumus ini, setiap variabel yang dideklarasikan pada segmen
data memiliki alamat fisik yang dapat dihitung oleh mahasiswa,
asalkan nilai \reg{DS} dan offset variabel tersebut diketahui.
Kedua nilai itu dapat dibaca langsung dari emulator atau
\eng{debugger}. Cara memperolehnya dibahas pada
Bagian~\ref{sec:contoh-program} dan Bagian~\ref{sec:praktikum-bab6}.
